\documentclass[12pt]{article}
\usepackage[T1]{fontenc}
\usepackage[utf8]{inputenc}
\usepackage[english]{babel}
\usepackage{lmodern}
\usepackage{microtype}
\usepackage{amsmath,amssymb,amsthm,mathtools}
\usepackage{mathrsfs}
\usepackage{enumitem}
\usepackage{geometry}
\geometry{margin=1.0in}
\usepackage{setspace}
\setstretch{1.06}
\usepackage{graphicx}
\usepackage[all,cmtip]{xy}
\usepackage{hyperref}
\hypersetup{colorlinks=true,linkcolor=blue,citecolor=blue,urlcolor=blue,hypertexnames=false}

\newcommand{\Aa}{\mathcal A}
\newcommand{\Bb}{\mathcal B}
\newcommand{\Cc}{\mathcal C}
\newcommand{\Dd}{\mathcal D}
\newcommand{\Ee}{\mathcal E}
\newcommand{\Ff}{\mathcal F}
\newcommand{\Gg}{\mathcal G}
\newcommand{\Hh}{\mathcal H}
\newcommand{\Ii}{\mathcal I}
\newcommand{\Jj}{\mathcal J}
\newcommand{\Kk}{\mathcal K}
\newcommand{\Ll}{\mathcal L}
\newcommand{\Mm}{\mathcal M}
\newcommand{\Oo}{\mathcal O}
\newcommand{\Pp}{\mathcal P}
\newcommand{\Qq}{\mathcal Q}
\newcommand{\Rr}{\mathcal R}
\newcommand{\Ss}{\mathcal S}
\newcommand{\Uu}{\mathcal U}
\newcommand{\Vv}{\mathcal V}
\newcommand{\Xx}{\mathcal X}
\newcommand{\Yy}{\mathcal Y}
\newcommand{\Zz}{\mathcal Z}
\newcommand{\ZZ}{\mathbb Z}
\newcommand{\QQ}{\mathbb Q}
\newcommand{\CC}{\mathbb C}
\newcommand{\GG}{\mathbb G}
\newcommand{\NN}{\mathbb N}
\newcommand{\LL}{\mathbb L}
\DeclareMathOperator{\Spec}{Spec}
\DeclareMathOperator{\Hom}{Hom}
\DeclareMathOperator{\Ext}{Ext}
\DeclareMathOperator{\Tor}{Tor}
\DeclareMathOperator{\Pic}{Pic}
\DeclareMathOperator{\Ob}{Ob}
\DeclareMathOperator{\codim}{codim}
\DeclareMathOperator{\Mod}{Mod}
\DeclareMathOperator{\Pro}{Pro}
\DeclareMathOperator{\Ind}{Ind}
\DeclareMathOperator{\Tr}{Tr}
\newcommand{\et}{\mathrm{\acute et}}
\newcommand{\red}{\mathrm{red}}
\newcommand{\op}{\mathrm{op}}
\newcommand{\id}{\mathrm{id}}
\newcommand{\rad}{\mathrm{rad}}
\newcommand{\RHom}{R\!\operatorname{Hom}}
\newcommand{\Hhom}{\mathcal{H}om}
\newcommand{\RHHom}{R\mathcal{H}om}
\newcommand{\RGamma}{R\Gamma}
\newcommand{\otimesL}{\mathbin{\overset{\mathbf L}{\otimes}}}
\newcommand{\boxtimesL}{\mathbin{\overset{\mathbf L}{\boxtimes}}}
\newcommand{\ie}{i.e.\ }
\newcommand{\cf}{cf.\ }
\newcommand{\resp}{respectively}
\newcommand{\SGA}{SGA}
\newcommand{\EGA}{EGA}
\newcommand{\Exp}{Expos\'e}
\newcommand{\muN}{\mu_n}
\newcommand{\qedtext}{\hfill\(\square\)}
\newenvironment{statement}[1]{\par\medskip\noindent\textbf{#1.}\ }{\par\medskip}
\newenvironment{remarktext}[1]{\par\medskip\noindent\textbf{#1.}\ }{\par\medskip}
\newenvironment{prooftext}[1][Proof]{\par\noindent\emph{#1.}\ }{\par\medskip}
\setlist[enumerate]{leftmargin=2.35em}
\setlist[itemize]{leftmargin=2.35em}
\allowdisplaybreaks
\emergencystretch=3em
\setlength{\parskip}{0.25em}
\setlength{\parindent}{1.15em}

\begin{document}

\begin{center}
{\Large \textbf{SGA 4 -- Working Draft}}\\[0.4em]
{\large \textbf{Expos\'e XVIII, Section 3, and Expos\'e XIX complete}}\\[0.25em]
{\normalsize Draft English LaTeX translation; compile-checked and rendered.}
\end{center}

\bigskip

\begin{center}
{\Large \textbf{SGA 4, Expos\'e XVIII -- Section 3}}\\[0.35em]
{\large \textbf{The Global Duality Theorem}}\\[0.25em]
{\normalsize P. Deligne, after A. Grothendieck}
\end{center}

\section*{3. The global duality theorem}

\subsection*{3.1. The functor \texorpdfstring{$Rf^!$}{Rf!}}

We shall need a variant of the canonical flasque resolutions.

\begin{statement}{Lemma 3.1.1}
Let $X$ be a coherent scheme, that is, a quasi-compact and quasi-separated scheme. The functor $\varinjlim$ induces an equivalence between the category of ind-objects of the category of constructible sheaves of sets (respectively of groups, respectively of abelian groups) and the category of sheaves of sets (respectively of ind-finite groups, respectively of torsion abelian groups).
\end{statement}

This follows from Expos\'e IX, Corollaries 2.7.2 and 2.7.3.

\paragraph{3.1.2. Modified canonical flasque resolutions.}
Let $\Ff$ be a torsion abelian sheaf on a scheme $X$, written as a filtered inductive limit of constructible sheaves $(\Ff_i)_{i\in I}$, and let $X_0$ be a conservative set of points of $X$. Define the \emph{modified canonical flasque resolution} of $\Ff$ by
\[
   C_\ell^*(\Ff)=\varinjlim_i C^*(\Ff_i),
\]
where $C^*(\Ff_i)$ denotes the canonical flasque resolution of Expos\'e XVII, 4.2.2. By Lemma 3.1.1, when $X$ is coherent this definition is independent of the chosen representation of $\Ff$ as a filtered colimit of constructible sheaves. It is also compatible with localization. This compatibility permits one to define $C_\ell^*(\Ff)$ by gluing, for any torsion sheaf $\Ff$ on any scheme $X$, not necessarily coherent.

The complex $C_\ell^*(\Ff)$ is a functorial resolution of $\Ff$. Its formation commutes with localization and with filtered inductive limits, and each functor $C_\ell^n$ is exact. Consequently, if $\Aa$ is a sheaf of rings on $X$, the functors $C_\ell^n$ carry $\Aa$-modules to $\Aa$-modules; they likewise carry torsion sheaves to torsion sheaves.

We shall use the following standard representability criterion.

\begin{statement}{Lemma 3.1.3}
Let $\Aa$ be a $\mathcal U$-abelian category having a small set of generators, and suppose that filtered inductive limits indexed by a small directed set are representable in $\Aa$ and are exact. A functor
\[
   F:\Aa^{\op}\longrightarrow (\mathrm{Sets})
\]
is representable if and only if it transforms arbitrary small inductive limits into projective limits.
\end{statement}

For instance, see SGA 4, Expos\'e I, 8.12.5, where one does not assume that $\Aa$ is abelian nor that filtered colimits are exact. We shall apply this to the category $\Mod(X,\Aa)$ of sheaves of modules on a ringed site $(X,\Aa)$. In this particular case, if $F$ is represented by the sheaf $\Ff$, then
\begin{equation}
   \Ff(U)=F(\Aa_U). \tag{3.1.3.1}
\end{equation}
This formula gives a direct verification of Lemma 3.1.3 in the case under consideration.

Recall that if $f:X\to S$ is compactifiable (Expos\'e XVII, 3.2), then $S$ is quasi-compact and quasi-separated, and the dimensions of the fibres of $f$ are bounded.

\begin{statement}{Theorem 3.1.4}
Let $f:X\to S$ be a compactifiable morphism and let $\Aa$ be a sheaf of torsion rings on $S$. Then the functor
\begin{equation}
   Rf_!:D(X,f^*\Aa)\longrightarrow D(S,\Aa) \tag{3.1.4.1}
\end{equation}
admits a partial right adjoint
\begin{equation}
   Rf^!:D^+(S,\Aa)\longrightarrow D^+(X,f^*\Aa). \tag{3.1.4.2}
\end{equation}
For $K\in D(X,f^*\Aa)$ and $L\in D^+(S,\Aa)$ one has a functorial isomorphism
\begin{equation}
   \Hom(Rf_!K,L)\simeq \Hom(K,Rf^!L). \tag{3.1.4.3}
\end{equation}
With the translation morphism defined by the usual compatibility square for shifts, this functor $Rf^!$ is triangulated.
\end{statement}

The notation $Rf^!$ is slightly abusive: in general it is not the derived functor of a pre-existing functor $f^!$.

\begin{prooftext}
Choose an integer $d$ such that every fibre of $f$ has dimension $<d$, and choose a compactification
\[
\xymatrix@R=1.1cm@C=1.4cm{
X\ar@{^{(}->}[r]^j\ar[d]_f&\bar X\ar[ld]^{\bar f}\\
S.}
\]
For a complex $K$ of $f^*\Aa$-modules on $X$, one has
\[
   Rf_!K\simeq R\bar f_*j_!K.
\]
For any sheaf $F$ on $X$, let $F^i$ be the components of the truncated modified canonical flasque resolution $\tau_{\leq 2d}C^*_\ell j_!F$. These components satisfy
\begin{equation}
   R^k\bar f_*F^i=0\qquad(k>0). \tag{3.1.4.5}
\end{equation}
Indeed, for $i\ne 2d$ the sheaf $F^i$ is a filtered colimit of flasque sheaves; for $i=2d$ and $k>0$ one has, by Expos\'e XVII, 5.2.8.1,
\[
   R^k\bar f_*F^{2d}=R^{k+2d}f_!F=0.
\]
The simple complex attached to the double complex obtained from the truncated modified canonical flasque resolution is a resolution of $K$. Hence, by (3.1.4.5),
\begin{equation}
   R\bar f_*j_!K \xleftarrow{\sim} \bar f_*\tau_{\leq 2d}''C_\ell^*j_!K. \tag{3.1.4.6}
\end{equation}
Define a functor from sheaves of modules on $X$ to complexes of sheaves on $S$ by
\begin{equation}
   f_!^{\bullet}(F)=\bar f_*\tau_{\leq 2d}C_\ell^*j_!F. \tag{3.1.4.7}
\end{equation}

\begin{statement}{Lemma 3.1.4.8}
The functors $f_!^i$ are exact and commute with filtered inductive limits; also $f_!^i=0$ for $i\notin [0,2d]$. Since $f_!^{\bullet}$ is a bounded complex of exact functors, it derives trivially to a functor $Rf_!^{\bullet}$, or simply $f_!^{\bullet}:D(X,f^*\Aa)\to D(S,\Aa)$, and (3.1.4.6) gives an isomorphism of functors
\[
   Rf_!^{\bullet}\xrightarrow{\sim}Rf_!.
\]
\end{statement}

The functors $C^i_\ell$ are exact, commute with filtered inductive limits, and carry every sheaf to an acyclic sheaf. The functors $\bar f_*$ and $j_!$ commute with filtered inductive limits. Thus $\bar f_*C^i_\ell j_!$, and hence $f_!^i$, commute with filtered inductive limits. For $i<2d$ the functor $f_!^i=\bar f_*C_\ell^ij_!$ is exact; for $i=2d$ exactness follows from (3.1.4.5). The remaining assertions are then immediate.

The definition of $f_!^{\bullet}$ is independent of the sheaf of rings $\Aa$ and still makes sense for any abelian sheaf. Lemma 3.1.4.8(i) remains true in the categories of torsion abelian sheaves on $X$ and on $S$.

By Lemmas 3.1.4.8 and 3.1.3, each exact functor
\[
   f_!^i:\Mod(X,f^*\Aa)\longrightarrow \Mod(S,\Aa)
\]
has a right adjoint $f_i^!$. Since $f_!^i$ is exact, $f_i^!$ carries injective sheaves to injective sheaves. With the sign convention of Expos\'e XVII, 1.1.12, the functors $f_i^!$ form a complex of left exact functors, zero in cohomological degrees outside $[-2d,0]$:
\begin{equation}
   f^!_{\bullet}:\Mod(S,\Aa)\longrightarrow C^b\Mod(X,f^*\Aa). \tag{3.1.4.10}
\end{equation}
This complex of functors derives to a triangulated functor
\begin{equation}
   Rf^!:D^+(S,\Aa)\longrightarrow D^+(X,f^*\Aa), \tag{3.1.4.11}
\end{equation}
right adjoint to $Rf_!$, which is the same as $Rf_!^{\bullet}$. In concrete terms, if $I\in C^+(S,\Aa)$ is a complex of injective sheaves, then all the sheaves $f_i^!I^n$ are injective, and one has an isomorphism of triple complexes
\begin{equation}
   \Hom^{\bullet}(f_!^{\bullet}K,I)\xrightarrow{\sim}
   \Hom^{\bullet}(K,f^!_{\bullet}I). \tag{3.1.4.12}
\end{equation}
Passing to $H^0$ of the associated simple complex gives the adjunction. The compatibility with shifts, hence the triangulated structure, is the standard verification. \qedtext
\end{prooftext}

\begin{remarktext}{Remark 3.1.5}
Suppose that $Rf^!$ has finite cohomological amplitude. Then, by Expos\'e XVII, 1.2.10, the functor $f^!_{\bullet}$ has a derived functor on the whole derived category,
\begin{equation}
   Rf^!:D(S,\Aa)\longrightarrow D(X,f^*\Aa). \tag{3.1.5.1}
\end{equation}
This derived functor is still right adjoint to $Rf_!$. Indeed, the usual description of morphisms in a derived category by replacing the source and target by quasi-isomorphic complexes gives
\begin{align*}
\Hom_{D(S)}(Rf_!K,L)&\simeq
 \varinjlim_{L\xrightarrow{\sim}L'}\Hom_{K(S)}(f_!^{\bullet}K,L')\\
&\simeq \varinjlim_{K'\xrightarrow{\sim}K,\;L\xrightarrow{\sim}L'}
\Hom_{K(S)}(f_!^{\bullet}K',L'),
\end{align*}
and similarly
\[
\Hom_{D(X)}(K,Rf^!L)
 \simeq \varinjlim_{K'\xrightarrow{\sim}K,\;L\xrightarrow{\sim}L'}
\Hom_{K(X)}(K',f^!_{\bullet}L').
\]
The two expressions are identified by the adjunction between $f_!^{\bullet}$ and $f^!_{\bullet}$.
\end{remarktext}

\begin{statement}{Definition 3.1.6}
The triangulated functor
\[
   Rf^!:D^+(S,\Aa)\longrightarrow D^+(X,f^*\Aa),
\]
or, under the hypotheses of 3.1.5, the functor
\[
   Rf^!:D(S,\Aa)\longrightarrow D(X,f^*\Aa),
\]
right adjoint to $Rf_!$, is called the \emph{exceptional inverse-image functor}.
\end{statement}

\begin{statement}{Proposition 3.1.7}
Suppose that $f:X\to S$ is compactifiable and of relative dimension at most $d$.
\begin{enumerate}[label=(\roman*)]
\item If $L\in D^+(S,\Aa)$ and $H^i(L)=0$ for $i\leq k$, then $H^i(Rf^!L)=0$ for $i\leq k-2d$.
\item If $L\in D^b(S,\Aa)$, then $\dim_{\mathrm{inj}}(Rf^!L)\leq \dim_{\mathrm{inj}}(L)$.
\end{enumerate}
\end{statement}

This follows from the fact that $Rf^!$ is the right derived functor of $f^!_{\bullet}$: the latter is zero in degrees $<-2d$ and in degrees $>0$, and the component functors $f_i^!$ send injectives to injectives.

\begin{statement}{Proposition 3.1.8}
Let $f:X\to S$ be a compactifiable quasi-finite morphism.
\begin{enumerate}[label=(\roman*)]
\item The functor $f_!:\Mod(X,f^*\Aa)\to \Mod(S,\Aa)$ has a right adjoint $f^!:\Mod(S,\Aa)\to \Mod(X,f^*\Aa)$, and $Rf^!$ is the derived functor of $f^!$.
\item If $f$ is a closed immersion, then $f^!$ is the functor of sections with support in $X$.
\item If $f$ is etale, then $f^!$ identifies with $f^*$, with trace
\[
   \Tr_f:f_!f^*F\longrightarrow F
\]
as adjunction morphism.
\end{enumerate}
\end{statement}

The first assertion is the case $d=0$ of the construction above, for which $f_!^{\bullet}=f_!$ and $f^!_{\bullet}=f^!$. Equivalently, it follows directly from Lemma 3.1.3 and the fact that $f_!$ is exact and commutes with filtered inductive limits. The second assertion is clear, and the third is Expos\'e XVII, 6.2.11.

\paragraph{3.1.9. The local form of adjunction.}
Let $f:X\to S$ be compactifiable, let $\Aa$ be a sheaf of torsion rings on $S$, and ring $X$ by $f^*\Aa$. Let $K\in D^-(X,f^*\Aa)$ and $L\in D^+(X,f^*\Aa)$. Represent $L$ by a bounded-below complex of injective sheaves, so that
\[
   R\Hhom(K,L)\simeq \Hhom(K,L).
\]
The complex $\Hhom(K,L)$ has flasque components. Thus
\begin{equation}
   Rf_*R\Hhom(K,L)\simeq f_*\Hhom(K,L). \tag{3.1.9.1}
\end{equation}
Choose, as in the proof of Theorem 3.1.4, a compactification and the corresponding functor $f_!^{\bullet}$. For every etale object $V$ over $S$, let $X_V=X\times_SV$, and write the induced compactification as
\[
\xymatrix@C=1.4cm@R=.55cm{
X_V\ar[rr]^{k_X}\ar@{^{(}->}[rd]_{j_V}\ar[ddd]^{f_V} && X\ar[ddd]^f\ar@{^{(}->}[rd]^j\\
&\bar X_V\ar[ldd]^{\bar f_V}\ar[rr]_{k_{\bar X}} && \bar X\ar[ldd]_{\bar f}\\
&&&\\
V\ar[rr]^k && S .
}
\]
There is a natural base-change isomorphism
\begin{equation}
   k^*f_!^{\bullet}\xrightarrow{\sim} f^{\bullet}_{V!}k_X^*. \tag{3.1.9.3}
\end{equation}
Using this, functoriality of $f^{\bullet}_{V!}$ gives a morphism of complexes
\begin{equation}
   f_*\Hhom^{\bullet}(K,L)\longrightarrow
   \Hhom^{\bullet}(f_!^{\bullet}K,f_!^{\bullet}L). \tag{3.1.9.4}
\end{equation}
By (3.1.9.1), it derives to
\begin{equation}
   Rf_*R\Hhom(K,L)\longrightarrow
   R\Hhom(Rf_!K,Rf_!L). \tag{3.1.9.5}
\end{equation}

Now let $K\in D^-(X,\Aa)$ and $L\in D^+(S,\Aa)$. Combining (3.1.9.4) with the adjunction morphism $Rf_!Rf^!L\to L$ gives the local form of the adjunction isomorphism:
\begin{equation}
   Rf_*R\Hhom(K,Rf^!L)\longrightarrow
   R\Hhom(Rf_!K,L). \tag{3.1.9.6}
\end{equation}

\begin{statement}{Proposition 3.1.10}
The morphism (3.1.9.6) is an isomorphism.
\end{statement}

Let $V$ be etale over $S$. The transitivity isomorphism
\[
   k_!Rf_{V!}\simeq Rf_!k_{X!}
\]
transposes to a localization isomorphism
\begin{equation}
   k_X^*Rf^!\simeq Rf_V^!k^*. \tag{3.1.10.1}
\end{equation}
To check (3.1.9.6), it is enough to check on all $V$ that the induced map
\begin{align}
\Hom(k_X^*K,k_X^*Rf^!L)
&\simeq H^0(V,Rf_*R\Hhom(K,Rf^!L)) \notag\\
&\longrightarrow H^0(V,R\Hhom(Rf_!K,L))
 \simeq \Hom(k^*Rf_!K,k^*L) \tag{3.1.10.2}
\end{align}
is an isomorphism. For $V=S$ this is precisely (3.1.4.3). The general case follows from the compatibility of (3.1.10.2) with the localization isomorphism (3.1.10.1) and the adjunction formula over $V$.

\begin{statement}{Lemma 3.1.10.3}
The square expressing this compatibility is commutative: the adjunction map after localizing by $k$ agrees with the adjunction map for $f_V$ under the localization isomorphism $k_X^*Rf^!\simeq Rf_V^!k^*$ and the base-change identification for $Rf_!$.
\end{statement}

\paragraph{3.1.11. Induction.}
Consider a Cartesian square
\begin{equation}
\xymatrix@C=1.8cm@R=1.3cm{
X'\ar[r]^{g'}\ar[d]^{f'}&X\ar[d]^f\\
S'\ar[r]^g&S,
}\tag{3.1.11.1}
\end{equation}
with $f$ compactifiable. Let $\Aa$ be a sheaf of torsion rings on $S$, $\Bb$ a sheaf of torsion rings on $S'$, and let $M$ be a bounded-above complex of $(f^*\Aa,\Bb)$-bimodules. At the level of complexes, there is a morphism, functorial in $K\in C(X,f^*\Aa)$,
\begin{equation}
   M\otimes_\Aa g^*f_!^{\bullet}K\longrightarrow
   f_{!}^{\prime\bullet}({f'}^*M\otimes_\Aa {g'}^*K). \tag{3.1.11.2}
\end{equation}
By transposition this gives a morphism, functorial in $L\in C(S',\Bb)$,
\begin{equation}
   g'_*\Hhom_\Bb({f'}^*M,f_{\bullet}^{\prime!}L)
   \longrightarrow
   f^!_{\bullet}g_*\Hhom_\Bb(M,L). \tag{3.1.11.3}
\end{equation}
Passing to derived functors gives a morphism
\begin{equation}
   Rg'_*R\Hhom_\Bb({f'}^*M,Rf'{}^!L)
   \longrightarrow
   Rf^!Rg_*R\Hom_\Bb(M,L). \tag{3.1.11.4}
\end{equation}
It is the morphism obtained from the base-change theorem for $Rf_!$ by adjunction. Indeed, the comparison diagram of Hom-groups contains, on one side, the map induced by (3.1.11.4), and on the other side the base-change isomorphism
\[
   Rf'_!({f'}^*M\otimesL_\Aa {g'}^*K)
      \xrightarrow{\sim} M\otimesL_\Aa g^*Rf_!K.
\]
Since this latter morphism is an isomorphism by the base-change theorem, the adjoint morphism (3.1.11.4) is also an isomorphism.

\begin{statement}{Proposition 3.1.12}
The morphism (3.1.11.4) is an isomorphism.
\end{statement}

This is called the \emph{induction isomorphism}. Three useful special cases are the following.

\begin{statement}{Corollary 3.1.12.1}
Let $f:X\to S$ be compactifiable, let $\Aa$ and $\Bb$ be sheaves of torsion rings on $S$, and let $\varphi:\Aa\to\Bb$ be a homomorphism. If $\rho$ denotes restriction of scalars from $\Bb$ to $\Aa$, then, for $K\in D^+(S,\Bb)$,
\[
   \rho Rf^!K\xrightarrow{\sim} Rf^!(\rho K).
\]
\end{statement}

Thus the sheaf of rings plays no essential role in the construction of $Rf^!$.

\begin{statement}{Corollary 3.1.12.2}
Let $f:X\to S$ be compactifiable, let $\Aa$ be a sheaf of torsion rings on $S$, and let $K\in D^-(S,\Aa)$, $L\in D^+(S,\Aa)$. Then there is the induction formula
\[
   R\Hhom(f^*K,Rf^!L)\xrightarrow{\sim}Rf^!R\Hhom(K,L).
\]
\end{statement}

\begin{statement}{Corollary 3.1.12.3}
For any square (3.1.11.1) and any $L\in D^+(S',\Aa)$,
\[
   Rg'_*R{f'}^!L\xrightarrow{\sim}Rf^!Rg_*L.
\]
\end{statement}

\paragraph{3.1.13. Composition.}
For a composite of compactifiable morphisms over a base $S$,
\[
   X\xrightarrow{h}Y\xrightarrow{g}Z,
\]
the composition isomorphism
\[
   Rg_!Rh_!\xrightarrow{\sim}R(gh)_!
\]
transposes into a composition isomorphism
\begin{equation}
   c^!_{g,h}:Rh^!Rg^!\xrightarrow{\sim}R(gh)^!. \tag{3.1.13.1}
\end{equation}
It satisfies the usual cocycle condition for a triple composite. Equivalently, the categories $D^+(X,f^*\Aa)$ are the fibres of a category both fibred and cofibred over the category of $S$-schemes with compactifiable $S$-morphisms; the direct-image and inverse-image functors are respectively $Rf_!$ and $Rf^!$.

For a commutative square of compactifiable morphisms
\[
\xymatrix@=1.6cm{
X'\ar[r]_g\ar[d]^{f'}&X\ar[d]^f\\
S'\ar[r]_{g'}&S,
}
\]
one obtains a cobase-change morphism
\begin{equation}
   Rf'_!Rg^!\longrightarrow Rg^!Rf_!. \tag{3.1.13.2}
\end{equation}

\paragraph{3.1.14. Base change for \texorpdfstring{$Rf^!$}{Rf!}.}
For a Cartesian square
\[
\xymatrix@=1.6cm{
X'\ar[r]^{g'}\ar[d]^{f'}&X\ar[d]^f\\
S'\ar[r]^g&S,
}
\]
with $f$ compactifiable, ring all four schemes by the inverse images of a sheaf of torsion rings $\Aa$ on $S$. The base-change isomorphism
\begin{equation}
   g^*Rf_!\xrightarrow{\sim}Rf'_!{g'}^* \tag{3.1.14.1}
\end{equation}
allows one to view the four derived categories as fibres of a cofibred category. Restricting to $D^+$ gives a category both fibred and cofibred, the inverse-image functors being $Rf^!$, $R{f'}^!$, $Rg_*$, and $Rg'_*$ in the appropriate directions. From this one recovers the transitivity isomorphism
\[
   Rg_*Rf^!\xrightarrow{\sim}Rf^!Rg_*,
\]
already obtained in Corollary 3.1.12.3, and also the base-change morphism
\begin{equation}
   {g'}^*Rf^!\longrightarrow R{f'}^!g^*. \tag{3.1.14.2}
\end{equation}
The morphism (3.1.14.2) may be described in two equivalent ways. By adjunction, it is the same as a morphism
\[
   Rf'_!{g'}^*Rf^!\longrightarrow g^*,
\]
and this is the composite
\[
   Rf'_!{g'}^*Rf^!\xleftarrow{\sim}g^*Rf_!Rf^!
      \xrightarrow{g^*(\mathrm{adj})} g^*.
\]
Equally, by adjunction it is the same as a morphism
\[
   Rf^!\longrightarrow Rg'_*R{f'}^!g^*,
\]
and this is obtained as
\[
   Rf^!\xrightarrow{Rf^!(\mathrm{adj})}Rf^!Rg_*g^*
      \xleftarrow{\sim}Rg'_*R{f'}^!g^*.
\]
Interpreted through the fibred/cofibred formalism, the isomorphism of Corollary 3.1.12.3 and the morphisms (3.1.13.2), (3.1.14.2) satisfy the two usual compatibilities of the type appearing in Expos\'e XII, 4.4, when either $f$ or $g$ is a composite of two morphisms.

\paragraph{3.1.15. A site attached to a morphism.}
Let $f:X\to S$ be a morphism of sites. Define a site $\Gamma(f)$ as follows. Its objects are triples
\[
   (U,V,\varphi),\qquad U\in\Ob(X),\ V\in\Ob(S),\ \varphi\in\Hom(U,f^*V).
\]
A family $(u_i,v_i):(U_i,V_i,\varphi_i)\to(U,V,\varphi)$ is covering if the maps $u_i:U_i\to U$ cover $U$ in $X$.

The functor $\Gamma(f)\to X$ sending $(U,V,\varphi)$ to $U$ is an equivalence of sites, written $\gamma:X\to\Gamma(f)$. The functor
\[
   f^*_{\Gamma}:S\longrightarrow\Gamma(f),\qquad
   V\longmapsto(f^*V,V,\id_{f^*V}),
\]
is a morphism of sites, denoted $f_\Gamma:\Gamma(f)\to S$, and one has $f_\Gamma\gamma\simeq f$. It admits a left adjoint on underlying categories,
\[
   f_{\Gamma\bullet}:(U,V,\varphi)\longmapsto V.
\]
This comorphism of sites defines the same morphism of topoi as the morphism of sites $f_\Gamma$.

Consequently, if $\Ff$ is a sheaf on $S$, if $f^\bullet_\Gamma\Ff$ is the presheaf on $\Gamma(f)$ defined by
\[
   f^\bullet_\Gamma\Ff(U,V,\varphi)=\Ff(V),
\]
and if $af^\bullet_\Gamma\Ff$ denotes the associated sheaf, then
\[
   f^*\Ff(U)\simeq af^\bullet_\Gamma\Ff(U,V,\varphi).
\]

\paragraph{3.1.16. A general Hom-construction.}
Let $f:X\to S$ be a morphism of sites, $\Aa$ a sheaf of rings on $S$, $\Bb$ a sheaf of rings on $X$, and $d$ an integer. Suppose that an object $K$ is given as follows. To every object $U$ of $X$, $K$ associates a complex $K(U)$ of $(\Aa,f_{|U*}(\Bb|U))$-modules on $S$, zero in degrees $>d$, and $K(U)$ is covariantly functorial in $U$.

For any sheaf $\Ff$ of $\Aa$-modules on $S$, let $K^!\Ff$ be the complex of sheaves on $X$ generated by the complex of presheaves of $\Bb$-modules
\[
   U\longmapsto \Hom_\Aa(K(U),\Ff).
\]
Let
\[
   RK^!:D^+(S,\Aa)\longrightarrow D^+(X,\Bb)
\]
be the derived functor. If $u:K_1\to K_2$ is a morphism of such objects, it induces transposed morphisms
\[
   {}^tu:K_2^!\longrightarrow K_1^!,\qquad RK_2^!\longrightarrow RK_1^!.
\]

\begin{statement}{Proposition 3.1.17}
Under these hypotheses:
\begin{enumerate}[label=(\roman*)]
\item For every point $x$ of $X$ and every $L\in D^+(S,\Aa)$, if $V(x)$ is the category of neighbourhoods of $x$, then
\begin{equation}
   \mathcal H^n(RK^!L)_x\simeq
   \varinjlim_{U\in\Ob V(x)}\Ext^n(K(U),L). \tag{3.1.17.1}
\end{equation}
\item Hence there is a spectral sequence
\begin{equation}
   E_2^{pq}=\varinjlim_{U\in\Ob V(x)}\Ext^p(\mathcal H^{-q}K(U),L)
   \Longrightarrow \mathcal H^n(K^!L)_x. \tag{3.1.17.2}
\end{equation}
\item If $X^0$ is a conservative set of points of $X$ and $u:K_1\to K_2$ is a morphism of such objects, then $^tu:RK_2^!\to RK_1^!$ is an isomorphism if and only if, for every $n$ and every $x\in X^0$, the morphism of pro-objects in $\Mod(S,\Aa)$
\begin{equation}
   u_x:\mathop{``\varprojlim\mbox{''}}_{U\in\Ob V(x)}\mathcal H^nK_1(U)
   \longrightarrow
   \mathop{``\varprojlim\mbox{''}}_{U\in\Ob V(x)}\mathcal H^nK_2(U) \tag{3.1.17.3}
\end{equation}
is an isomorphism.
\end{enumerate}
\end{statement}

For (i), take $L$ to be a complex of injectives; then (3.1.17.1) is immediate. Assertion (ii) follows from (i). The sufficiency in (iii) follows at once from (3.1.17.2). The necessity, although not needed later, follows by testing against injectives $I$: if $^tu$ is an isomorphism, then for every injective $I$ the Hom functors out of the two pro-objects in (3.1.17.3) agree. Since such Hom functors form a conservative system on the pro-category, the pro-objects are isomorphic.

\begin{remarktext}{Example 3.1.18}
Let $f:X\to S$ be a compactifiable morphism of schemes and $\Aa$ a sheaf of torsion rings on $S$. Ring $X$ by $f^*\Aa$. Choose an integer $d$ bounding the dimensions of the fibres of $f$ and choose a compactification. Put
\begin{equation}
   K(U)=f_!^{\bullet}((f^*\Aa)_U). \tag{3.1.18.1}
\end{equation}
Then, by the formula (3.1.3.1) for the adjoint of $f_!^{\bullet}$,
\begin{equation}
   f^!_{\bullet}\simeq K^!: \Mod(S,\Aa)\longrightarrow C^b\Mod(X,f^*\Aa), \tag{3.1.18.2}
\end{equation}
and therefore
\begin{equation}
   Rf^!=RK^!:D^+(S,\Aa)\longrightarrow D^+(X,f^*\Aa). \tag{3.1.18.3}
\end{equation}
\end{remarktext}

\begin{remarktext}{Example 3.1.19}
Let $f:X\to S$, $\Aa$, $\Bb$, and $d$ be as in 3.1.16. For each $W=(U,V,\varphi)\in\Ob\Gamma(f)$, let $K(W)$ be a complex of $(\Aa_V,\varphi_*(\Bb|U))$-modules on $V$, functorial in $W$ in the sense that for a morphism $(u,v):W_1\to W_2$ one is given maps
\[
   K(W_1)\longrightarrow v^*K(W_2),
   \qquad\text{equivalently}\qquad
   v_!K(W_1)\longrightarrow K(W_2).
\]
To such a $K$ is associated an object $K'$ of the type 3.1.16.1 relative to $f_\Gamma:\Gamma(f)\to S$, namely
\begin{equation}
   K'(U,V,\varphi)=j_!K(U,V,\varphi), \tag{3.1.19.2}
\end{equation}
where $j$ is the inclusion morphism of $V$ into $S$. One sets $K^!={K'}^!$ and $RK^!=R{K'}^!$.
\end{remarktext}

\begin{remarktext}{Example 3.1.20}
In the notation of 3.1.19, take $\Bb=f^*\Aa$ and put
\begin{equation}
   K(U,V,\varphi)=\Aa_V\quad\text{in degree }0. \tag{3.1.20.1}
\end{equation}
Then 3.1.15.4 gives functorial isomorphisms
\begin{equation}
   K^!\simeq f^*,\qquad RK^!\simeq f^*. \tag{3.1.20.2--3}
\end{equation}
\end{remarktext}

\subsection*{3.2. Poincare duality}

\paragraph{3.2.1. Construction of the comparison morphism.}
Let $f$ be compactifiable and flat, purely of relative dimension $d$; more generally, it is enough for $f$ to satisfy condition $(*)_d$ of Theorem 2.9. Let $n\geq 1$ be invertible on $S$ and ring $S$ by a sheaf of rings $\Aa$ such that $n\Aa=0$. By Example 3.1.18, the functor $Rf^!:D^+(S,\Aa)\to D^+(X,f^*\Aa)$ may be computed by the construction of 3.1.16.

Let $K'$ be the object of the type 3.1.19.1 that assigns to $(U,V,\varphi)\in\Gamma(f)$ the complex $\ZZ/n(-d)[-2d]$ on $V$, reduced to $\ZZ/n(-d)$ placed in degree $2d$, and let $K''$ be the associated object of the type 3.1.16.1 relative to $f:\Gamma(f)\to S$. The trace morphism of Theorem 2.9 gives, for each $(U,V,\varphi)$, a morphism of sheaves on $V$
\[
   \Tr:R^{2d}\varphi_!\ZZ/n\longrightarrow \ZZ/n(-d),
\]
and hence a morphism of functors
\begin{align}
K(U)&=\bar f_*\tau_{\leq 2d}C_\ell^*\ZZ/n_U
     \longrightarrow \mathcal H^{2d}K(U)[-2d]                                      \notag\\
&\simeq R^{2d}f_!(\ZZ/n)_U[-2d]
  \simeq j_!R^{2d}\varphi_!\ZZ/n[-2d]
  \xrightarrow{\Tr_\varphi} j_!\ZZ/n(-d)[-2d]=K''(U,V,\varphi). \tag{3.2.1.1}
\end{align}
A twisted form of Example 3.1.20 gives
\[
   R{K''}^!\simeq Rf^*(d)[2d].
\]
Thus 3.1.16 gives a composite morphism
\begin{equation}
   t_f: Rf^*(d)[2d]\simeq R{K''}^!\longrightarrow RK^!\simeq Rf^!. \tag{3.2.1.2}
\end{equation}
If $Rf^!$ has finite cohomological dimension, the same morphism is defined as a morphism of functors on the whole derived category.

There is another construction of (3.2.1.2): by adjunction from the trace morphism
\[
   \Tr:Rf_!(f^*L(d)[2d])\longrightarrow L.
\]
The two constructions agree.

\begin{statement}{Lemma 3.2.3}
For $L\in D^+(S,\Aa)$, and also for $L\in D(S,\Aa)$ if $Rf^!$ has finite cohomological dimension, the square
\[
\xymatrix@C=2.3cm@R=1.55cm{
Rf_!(f^*L(d)[2d])\ar[d]_{\Tr_f}\ar[r]^{Rf_!(t_f)}&Rf_!Rf^!L\ar[d]^{\mathrm{adjunction}}\\
L\ar@{=}[r]&L
}
\]
is commutative.
\end{statement}

A concrete way to see this is to write $t_f$ in terms of the functor $f^!_{\bullet}$. For $U\in X_{\et}$, $V\in S_{\et}$, and $\varphi\in\Hom_f(U,V)$,
\[
   f^!_{\bullet}(L)(U)=\Hom(f_!^{\bullet}\Aa_U,L).
\]
The trace morphism gives
\[
   \Tr_f:f_!^{\bullet}\Aa_U\longrightarrow \Aa_V(-d)[-2d],
\]
which transposes into
\[
   t_\varphi:f^!_{\bullet}(L)(U)\longleftarrow L(d)[2d](V).
\]
The morphisms $t_\varphi$ define
\[
   t_f:f^*L(d)[2d]\longrightarrow f^!_{\bullet}L,
\]
which coincides in the derived category with (3.2.1.2). In this description the relevant diagrams of complexes are visibly commutative: first for $\Aa_V$, then by functoriality for an arbitrary sheaf or complex $K$. Taking $K=L$ and the identity of $L$ gives the assertion.

\paragraph{3.2.4. Compatibility with composition.}
From Lemma 3.2.3 and from the definition of the composition isomorphisms by adjunction, one obtains the following. If $f=gh$ is a composite of two compactifiable flat morphisms of finite presentation, purely of relative dimensions $d_1$ and $d_2$, and if $d=d_1+d_2$, then the square
\[
\xymatrix@C=2.5cm@R=1.5cm{
h^*(g^*L(d_2)[2d_2])(d_1)[2d_1]\ar[r]^-{\sim}\ar[d]_{t_h\circ t_g}&
(gh)^*L(d)[2d]\ar[d]^{t_{gh}}\\
Rh^!Rg^!L\ar[r]^-{\sim}&R(gh)^!L
}
\]
is commutative.

\begin{statement}{Theorem 3.2.5 (Poincare duality)}
Let $f:X\to S$ be a compactifiable smooth morphism, so that $S$ is quasi-compact and quasi-separated. Let $d$ be the locally constant function on $X$ giving the relative dimension of $f$. Let $n\geq 1$ be invertible on $S$, and let $\Aa$ be a sheaf of rings on $S$ with $n\Aa=0$. Then
\[
   Rf_!:D(X,f^*\Aa)\longrightarrow D(S,\Aa)
\]
has as right adjoint the functor
\[
   K\longmapsto f^*K(d)[2d]
\]
after the canonical identification with $Rf^!$. The adjunction morphism is the trace
\[
   \Tr_f:Rf_!f^*K(d)[2d]\longrightarrow K,
\]
and consequently
\[
   \Hom(K,f^*L(d)[2d])\xrightarrow{\sim}\Hom(Rf_!K,L).
\]
\end{statement}

\begin{prooftext}
One reduces to the case where $d$ is constant. By the construction of $t_f$ in 3.2.1, Theorem 2.14 and the criterion 3.1.17(iii) imply that
\[
   t_f:f^*L(d)[2d]\xrightarrow{\sim}Rf^!L
\]
is an isomorphism for $L\in D^+(S,\Aa)$. In particular, $Rf^!$ has finite cohomological dimension, hence is defined on the whole derived category, and the same $t_f$ remains an isomorphism there. By definition $Rf^!$ is right adjoint to $Rf_!$. Lemma 3.2.3 identifies the adjunction morphism $Rf_!Rf^!K\to K$, via $t_f$, with the trace $\Tr_f$. This proves the theorem. \qedtext
\end{prooftext}

From now on one identifies the functors $f^*L(d)[2d]$ and $Rf^!L$ by means of $t_f$.

\paragraph{3.2.6. Classical form over an algebraically closed field.}
Let $S=\Spec k$, with $k$ algebraically closed, and let $X$ be compactifiable and smooth over $k$, purely of relative dimension $d$. Let $n$ be invertible in $k$, and let $F$ be a constructible locally constant sheaf of $\ZZ/n$-modules, that is, a coefficient system killed by $n$. If $^\vee$ denotes the $\ZZ/n$-dual, then
\[
   R\Hhom(F,\ZZ/n)=F^\vee
\]
in degree $0$, because $\ZZ/n$ is injective as a $\ZZ/n$-module. The isomorphism deduced from 3.2.5 gives
\begin{equation}
   R\Gamma(F^\vee(d)[2d])\xrightarrow{\sim}R\Hom(R\Gamma_c(F),\ZZ/n). \tag{3.2.6.1}
\end{equation}
Passing to cohomology and using again the injectivity of $\ZZ/n$ gives
\begin{equation}
   H^{2d-i}(X,F^\vee(d))\xrightarrow{\sim}H_c^i(X,F)^\vee, \tag{3.2.6.2}
\end{equation}
which is the usual form of Poincare duality.

\subsection*{Bibliography of Expos\'e XVIII}

\begin{itemize}
\item J. Giraud, \emph{Non-abelian cohomology}.
\item A. Grothendieck, \emph{The Brauer group III}, in \emph{Dix expos\'es sur la cohomologie des sch\'emas}.
\item L. Illusie, \emph{Complexe cotangent et d\'eformations II}, Lecture Notes in Mathematics 283.
\item D. Mumford, \emph{Geometric invariant theory}, Springer-Verlag, 1965.
\item M. Raynaud, \emph{Faisceaux amples sur les sch\'emas en groupes et les espaces homog\`enes}, Lecture Notes in Mathematics 119.
\item J.-P. Serre, \emph{Groupes alg\'ebriques et corps de classe}, Hermann, 1959.
\item J.-L. Verdier, \emph{Duality in the cohomology of locally compact spaces}, S\'eminaire Bourbaki 300.
\item J.-L. Verdier, \emph{A duality theorem in the etale cohomology of schemes}, Woods Hole seminar on algebraic geometry, 1964.
\item J.-L. Verdier, \emph{Derived categories}, IHES mimeographed notes.
\end{itemize}

\newpage

\begin{center}
{\Large \textbf{SGA 4, Expos\'e XIX}}\\[0.35em]
{\large \textbf{Cohomology of Excellent Preschemes of Equal Characteristic}}\\[0.25em]
{\normalsize M. Artin}
\end{center}

\bigskip

In this expos\'e one develops etale cohomology for excellent preschemes in equal characteristic, assuming the resolution of singularities theorem for such preschemes in the following form.

Let $X=\Spec A$ be an excellent local scheme of equal characteristic, and let $U\subset X$ be a regular open subset. There exist a proper birational morphism $f:X'\to X$ and an open immersion $U\to X'$ over $X$ such that $X'$ is regular and $Y=X'-U$ is everywhere of codimension $1$ with normal crossings.

Thus the proofs are, for the moment, valid only in characteristic zero. One can also draw some consequences in characteristic $p>0$ in low dimensions. Under the resolution hypothesis, one obtains reasonably satisfactory results for preschemes of equal characteristic and for coefficients prime to the characteristic, especially the purity, local acyclicity, finiteness, and cohomological-dimension theorems stated below. By contrast, very little is known in general about the cohomology of arbitrary preschemes of unequal characteristic, even in dimension $2$, where resolution is available. For example, the purity theorem below is not known for complete local rings of unequal characteristic of dimension $>1$. The result for formal power-series rings $k[[x_1,\ldots,x_n]]$ is one of the main tools. Its proof does not use resolution; consequently that theorem is established in arbitrary characteristic.

We shall use without repeated mention the standard properties of excellent schemes from EGA IV 7.8, such as stability under finite type extension and under strict localization. The latter fact follows from resolution and EGA IV 7.9.5.

\section*{1. Purity for \texorpdfstring{$k[[x_1,\ldots,x_n]]$}{k[[x1,...,xn]]}}

We recall the following consequence of Expos\'e XVI, 3.7.

\begin{statement}{Corollary 1.1}
Let $A$ be a strictly local ring and let $A\{x\}$ denote the strict localization of $A[x]$ at the prime ideal generated by $\rad A$ and $x$. Put
\[
   U=\Spec A\{x\}[1/x].
\]
For $n$ prime to the residual characteristic of $A$ one has
\[
H^q(U,\mu_n)\simeq
\begin{cases}
\mu_n(A),&q=0,\\
\ZZ/n,&q=1,\\
0,&q>1.
\end{cases}
\]
\end{statement}

Indeed, let $X=\Spec A[x]$ and let $Y=V(x)$. Then $(Y,X)$ is a smooth pair over $\Spec A$. By the results of Expos\'es V and VIII, $H^q(U,\mu_n)$ is identified with the fibre of the sheaf $\mathcal H_Y^{q+1}(X,\mu_n)$ at a geometric point over $(\rad A,x)$; the corollary follows at once from Expos\'e XVI, 3.7. We have written both $\mu_n$ and $\ZZ/n$ in order to display the canonical formulae used later.

\begin{statement}{Theorem 1.2}
Let $k$ be a separably closed field and let
\[
   U=\Spec k[[x_1,\ldots,x_r]][1/x_r].
\]
For $n$ prime to the characteristic of $k$ one has
\[
H^q(U,\mu_n)=
\begin{cases}
\mu_n(A),&q=0,\\
\ZZ/n,&q=1,\\
0,&q>1.
\end{cases}
\]
\end{statement}

The same proof also applies to the ring of convergent power series when $k=\CC$. The hypothesis that $k$ be separably closed can be avoided by stating the result more carefully. For the proof of the general purity theorem below, however, only the separably closed case is needed; the other assertions follow from the general theorem.

For dimensions $0$ and $1$ one has the following result for arbitrary regular strictly local rings.

\begin{statement}{Lemma 1.3}
Let $A$ be a regular strictly local ring and let $x\in \rad A$ be a local parameter. Put $U=\Spec A[1/x]$. Then
\[
   H^0(U,\mu_n)=\mu_n(A),
\]
and there are canonical isomorphisms
\[
   \ZZ/n\simeq H^0(U,\GG_m)/(H^0(U,\GG_m))^n
       \simeq H^1(U,\mu_n),
\]
where the generator of $\ZZ/n$ is the residue of $x\in H^0(U,\GG_m)$.
\end{statement}

\begin{prooftext}
For $q=0$ the assertion is that $U$ is connected and non-empty. For $q=1$, apply the Kummer exact sequence. Since $A$ is regular and factorial, and since $U$ is an open subscheme of $X=\Spec A$, one has $\Pic U=0$. Thus
\[
   H^0(U,\GG_m)\xrightarrow{n}H^0(U,\GG_m)	o H^1(U,\mu_n)	o 0
\]
is exact. There is also the evident exact sequence
\[
0\to H^0(X,\GG_m)\to H^0(U,\GG_m)\to \ZZ\to 0,
\]
where the last map is the order of vanishing along $x=0$. Since $A$ is strictly local and $n$ is invertible in the residue field, $A^*$ is divisible by $n$; the equation $Y^n-u=0$ defines a finite etale algebra that is completely split. Hence the quotient is $\ZZ/n$, generated by the residue of $x$. \qedtext
\end{prooftext}

\begin{prooftext}[Proof of Theorem 1.2]
Only the vanishing for $q>1$ remains. Using the terminology of Expos\'e V, Appendix, write
\[
   H^q(U,\mu_n)=\varinjlim_{\Uu}H^q(\Uu,\mu_n),
\]
where $\Uu$ runs over the filtering category, up to homotopy, of hypercoverings of $U$. It is enough, given a hypercovering $\Uu$, to find a refinement $\Vv\to\Uu$ for which the induced map
\[
   H^q(\Uu,\mu_n)\longrightarrow H^q(\Vv,\mu_n)
\]
is zero.

Take the terms of $\Uu$ separated and of finite type. Let $W_1,\ldots,W_m$ be the connected components appearing in finitely many low degrees of the hypercovering. Each $W_j$ is finite over an open of the spectrum of a finite $A$-algebra. More explicitly, after replacing by suitable normalizations, one obtains finite $A$-algebras $B_j$ and $C_j$ presenting the relevant objects and morphisms. The key point is that these data descend from the formal ring
\[
   A=k[[x_1,\ldots,x_r]]
\]
to
\[
   A^0=k[[x_1,\ldots,x_{r-1}]]\{x_r\},
\]
the henselization of the polynomial extension at the origin. The descent uses the Weierstrass preparation theorem: if $B_j$ is defined by a monic polynomial $f_j$, then the ideal generated by the discriminant or different of $f_j$ descends to $A^0$, and one can choose a monic polynomial $f^0_j\in A^0[Y]$ with
\[
   f_j\equiv f^0_j\pmod{d_j^2\rad A}.
\]
The corresponding algebras $A[Y]/(f_j)$ and $A[Y]/(f^0_j)$ are isomorphic; after normalization this descends $B_j$ and the relevant opens. Consequently the chosen part of the hypercovering is induced from a hypercovering $\Uu^0$ over
\[
   U^0=\Spec A^0[1/x_r].
\]

But $A^0$ is a strict localization of a polynomial algebra over $k[[x_1,\ldots,x_{r-1}]]$, and the result follows from Corollary 1.1 together with induction on $r$. The refinement coming from $\Uu^0$ kills the chosen cohomology class. This proves $H^q(U,\mu_n)=0$ for $q>1$. \qedtext
\end{prooftext}

\begin{statement}{Lemma 1.5}
In the algebraization step used above, the finite algebras and the finite morphisms appearing in a finite part of the hypercovering over $A=k[[x_1,\ldots,x_r]]$ are induced from the corresponding henselian ring $A^0=k[[x_1,\ldots,x_{r-1}]]\{x_r\}$.
\end{statement}

The proof is the Weierstrass argument just described. If $C_j$ is already finite over $k[[x_1,\ldots,x_{r-1}]]$, take it as an $A^0$-algebra by restriction of scalars. For $B_j$, choose a monic polynomial $f_j$ and descend the ideal generated by its discriminant to an ideal $I_j^0\subset A^0$. Since the $I_j^0$-adic completion of $A^0$ is $A$, choose $f_j^0\in A^0[Y]$ congruent to $f_j$ modulo $d_j^2\rad A$. The comparison theorem for such monic equations gives
\[
   A[Y]/(f_j)\simeq A[Y]/(f_j^0),
\]
and the normalizations agree after base change to $A$. This gives the asserted descent.

\section*{2. The case of a strictly local ring}

\begin{statement}{Theorem 2.1}
Assuming resolution of singularities as in the introduction, let $A$ be an excellent, regular, strictly local ring of equal characteristic, and let $x$ be a local parameter, meaning an element that is part of a regular system of parameters. Put $U=\Spec A[1/x]$. For $n$ prime to the characteristic of $A$,
\[
H^q(U,\mu_n)=
\begin{cases}
\mu_n(A),&q=0,\\
\ZZ/n,&q=1,\\
0,&q>1.
\end{cases}
\]
\end{statement}

The proof is by induction on $q$. Let $P(N)$ be the assertion of the theorem for all $q\leq N$. Lemma 1.3 gives $P(1)$.

\begin{statement}{Lemma 2.2}
Assume $P(N)$. Let $A\to A'$ be a local homomorphism of excellent, regular, strictly local rings of equal characteristic. Let $x_1,\ldots,x_r$ be part of a regular system of parameters of $A$, and suppose that their images are still part of a regular system of parameters of $A'$. Put
\[
   U=\Spec A[1/x_1,\ldots,1/x_r],\qquad
   U'=\Spec A'[1/x'_1,\ldots,1/x'_r].
\]
Then
\[
   H^q(U,\mu_n)\longrightarrow H^q(U',\mu_n)
\]
is bijective for every $q\leq N$.
\end{statement}

\begin{prooftext}
Argue by induction on $r$. The case $r=0$ is clear. For $r>0$, set $B=A/(x_r)$,
\[
   V=\Spec A[1/x_1,\ldots,1/x_{r-1}],\qquad
   Y=\Spec B[1/x_1,
\ldots,1/x_{r-1}],
\]
and define $V'$ and $Y'$ analogously. There is an open immersion $i:U\hookrightarrow V$, with closed complement $Y$ defined by $x_r=0$.

Applying $P(N)$ to the strict localizations of $V$ at geometric points of $Y$ gives
\begin{equation}
R^qi_*\mu_n=
\begin{cases}
\mu_{n,V},&q=0,\\
(\ZZ/n)_Y,&q=1,\\
0,&1<q\leq N.
\end{cases} \tag{2.3}
\end{equation}
The same formula holds after prime. The comparison map of Leray spectral sequences for $i$ and $i'$ has $E_2$-terms
\[
   E_2^{pq}=H^p(V,R^qi_*\mu_n),\qquad {E'}_2^{pq}=H^p(V',R^qi'_*
\mu_n).
\]
The induction hypothesis on $r-1$, applied to $V$ and to $Y$, shows that the maps on the non-zero rows $q=0,1$ are isomorphisms in degrees $\leq N$. Since the rows $1<q\leq N$ vanish, the corresponding exact sequences give injectivity of
\[
   H^m(U,\mu_n)\to H^m(U',\mu_n)
\]
for $m\leq N$.

It remains to prove surjectivity. The boundary map
\[
   H^m(U,\mu_n)\longrightarrow H^{m-1}(Y,\ZZ/n)
\]
is surjective for $m\leq N$. To see this one reduces, by applying the already proved injectivity to a comparison with the henselized polynomial ring, to the case $A=k\{x_1,\ldots,x_r\}$. For this ring, which is a limit of schemes smooth over $k$, the purity theorem of Expos\'e XVI, 3.7 gives $R^qi_*\mu_n=0$ for $q>1$. Thus the boundary is surjective exactly when
\[
   H^{m+1}(V,\mu_n)\longrightarrow H^{m+1}(U,\mu_n)
\]
is injective. But $V$ is the inverse image of
\[
   Z=\Spec k\{x_1,\ldots,x_{r-1}\}[1/x_1,
\ldots,1/x_{r-1}]
\]
under a smooth locally acyclic morphism, so $H^q(Z,\mu_n)\simeq H^q(V,\mu_n)$. Moreover $U$ is a projective limit of opens $U_\alpha$ finite etale over the corresponding polynomial situation, and each $U_\alpha$ admits a section over $Z$. Therefore $H^q(Z,\mu_n)\to H^q(U_\alpha,\mu_n)$, and hence $H^q(Z,\mu_n)\to H^q(U,\mu_n)$, is injective. This proves the lemma. \qedtext
\end{prooftext}

\begin{statement}{Lemma 2.4}
Assume $P(N)$. Let $A\to A'$ be a regular morphism of excellent strictly local rings of equal characteristic, not necessarily regular, and let $U\subset\Spec A$ be a regular open subscheme. Put $U'=U\times_{\Spec A}\Spec A'$. For $n$ prime to the residual characteristic of $A$, the morphism
\[
   H^q(U,\mu_n)\longrightarrow H^q(U',\mu_n)
\]
is bijective for $q\leq N$.
\end{statement}

\begin{prooftext}
Choose, by resolution, a proper morphism $f:X\to\Spec A$ resolving $\Spec A$, with $U$ embedded as an open $i:U\hookrightarrow X$, and with $Y=X-U$ purely of codimension $1$ and with normal crossings. After base change by $A\to A'$ the corresponding objects $X'$, $U'$, $Y'$ have the same properties, because regular morphisms preserve regularity and normal-crossing local equations.

At a point $p'\in X'$ mapping to $p\in X$, choose local equations $x_1,
\ldots,x_r$ for $Y$ that are part of a regular system of parameters. Their images in $\Oo_{X',p'}$ are again part of a regular system of parameters and cut out $Y'$. Passing to strict localizations $B$ and $B'$ at corresponding geometric points, Lemma 2.2 applies to
\[
   \Spec B[1/x_1,\ldots,1/x_r]	o
   \Spec B'[1/x'_1,\ldots,1/x'_r]
\]
and gives isomorphisms on $H^q(-,\mu_n)$ for $q\leq N$. Hence the base-change morphisms
\[
   g^*R^qi_*\mu_n\longrightarrow R^qi'_*
\mu_n
\]
are isomorphisms for $q\leq N$. Proper base change, applied to $X$ and $X'$, gives isomorphisms
\[
   H^p(X,F)\xrightarrow{\sim}H^p(X',g^*F)
\]
for torsion sheaves $F$. Comparing the Leray spectral sequences for $i$ and $i'$ then gives the desired isomorphism on $H^m$ for $m\leq N$. \qedtext
\end{prooftext}

\begin{statement}{Lemma 2.5}
One has $P(N)\Rightarrow P(N+1)$.
\end{statement}

\begin{prooftext}
Let $A$ be strictly local, regular, excellent, and of equal characteristic; let $x\in\rad A$ be a local parameter; put $U=\Spec A[1/x]$. Let $\widehat A$ be the completion, $\widehat U=\Spec\widehat A[1/x]$, and let $g:\widehat U\to U$ be the canonical morphism. It suffices to show
\begin{equation}
   g_*\mu_{n,\widehat U}\simeq \mu_{n,U},\qquad
   R^qg_*\mu_n=0\quad (1\le q\le N). \tag{2.6}
\end{equation}
Then the Leray spectral sequence gives isomorphisms $H^q(U,\mu_n)\simeq H^q(\widehat U,\mu_n)$ for $q\leq N$ and an injection in degree $N+1$. Since $P(1)$ is known, it remains only to prove vanishing in the range $1<q\leq N+1$. Replacing $A$ by its completion reduces to
\[
   \widehat A\simeq k[[x_1,\ldots,x_r]],\qquad x=x_r,
\]
with $k$ separably closed, which is Theorem 1.2.

To verify (2.6), recall that $R^qg_*\mu_n$ is the sheaf associated to the presheaf
\[
   V\longmapsto H^q(\widehat V,\mu_n)
\]
for etale $V\to U$. Taking $V$ separated, of finite type, and integral, let $B$ be the normalization of $A$ in the fraction field of $\Gamma(V,\Oo_V)$. Then $B$ is finite over $A$, $V$ is a regular open subscheme of $\Spec B$, and $\widehat V$ is the corresponding open in $\Spec\widehat B$, with $\widehat B=B\otimes_A\widehat A$. Lemma 2.4 applies to $B\to\widehat B$ and gives
\[
   H^q(V,\mu_n)\simeq H^q(\widehat V,\mu_n)
\]
for $q\leq N$. But the sheaf associated to $V\mapsto H^q(V,\mu_n)$ is zero in positive degree on the etale site. This proves (2.6), and hence the lemma. \qedtext
\end{prooftext}

Theorem 2.1 follows from Lemmas 2.5 and 1.3 by induction on $N$.

\section*{3. Purity}

We use terminology parallel to that of Expos\'e XVI, Section 3.

\begin{statement}{Definition 3.1}
A \emph{regular pair $(Y,X)$ of codimension $c$} is a closed immersion $i:Y\to X$ of locally noetherian regular schemes such that, for every $y\in Y$, one has $\codim_y(Y,X)=c$. Equivalently, locally on $X$ near $y$, the ideal of $Y$ is generated by sections $x_1,\ldots,x_c$ forming part of a regular system of parameters at every point of $Y$. We put $U=X-Y$ and denote by $j:U\to X$ the corresponding open immersion.
\end{statement}

A morphism $(Y',X')\to(Y,X)$ of regular pairs of codimension $c$ is a morphism $X'\to X$ such that $Y'=Y\times_XX'$. The notions of regular triples and of morphisms of regular triples are defined in the evident analogous way.

\begin{statement}{Theorem 3.2 (Cohomological purity)}
Assume resolution of singularities as above. Let $(Y,X)$ be a regular pair of codimension $c>0$, with $X$ an excellent prescheme of equal characteristic. Let $F$ be a locally constant sheaf on $X$ of cyclic groups of order $n$, where $n$ is prime to the characteristic of $X$. Then, in the notation of Expos\'es V and VIII,
\begin{align*}
\mathcal H^q_Y(X,F)&=(R^qi^!)F=0\qquad(q\ne 2c),\\
(R^qj_*)j^*F&=0\qquad(q\ne0,\ 2c-1),
\end{align*}
and
\[
   \mathcal H_Y^{2c}(X,F)=i^*(R^{2c-1}j_*)j^*F
\]
is a locally constant sheaf on $Y$ of cyclic groups of order $n$.
\end{statement}

\begin{prooftext}
Clearly $\mathcal H_Y^q(X,F)=0$ for $q=0,1$. By the long exact cohomology sequence for supports,
\[
   \mathcal H_Y^q(X,F)=(R^qi^!)F\simeq i^*(R^{q-1}j_*)j^*F\qquad(q>1),
\]
so the assertions are equivalent. They are etale-local on $X$, so we reduce to $F=\mu_n$.

For $c=1$, computing fibres of $(R^qj_*)j^*\mu_n$ by strict localization reduces exactly to Theorem 2.1. Thus the stalk is $\ZZ/n$ for $q=1$ and zero for $q>1$. Since the sheaves are zero away from $Y$, this gives all assertions for $q>1$; for $q=1$ one applies the local constancy criterion of Expos\'e IX, 2.13, together with the explicit calculation in Lemma 1.3.

For $c>1$ one argues by induction. Locally on $X$ choose a regular triple $(Y,Z,X)$ such that $(Z,X)$ has codimension $1$ and $(Y,Z)$ has codimension $c-1$. If
\[
\xymatrix{Y\ar[rr]^u\ar[rd]_i&&Z\ar[ld]^v\\&X&}
\]
are the closed immersions, then the spectral sequence
\[
   E_2^{pq}=(R^pu^!)(R^qv^!)F\Longrightarrow(R^{p+q}i^!)F
\]
and the induction hypothesis show that the only nonzero term occurs at $(p,q)=(2c-2,2)$, where it is locally constant cyclic of order $n$. This proves the theorem. \qedtext
\end{prooftext}

\paragraph{3.3. Canonical determination of the pure sheaf.}
It remains to determine canonically the locally constant sheaf occurring in Theorem 3.2. For every prescheme $S$ define
\[
   (\mu_n^{\otimes r})_S
\]
to be the $r$-fold tensor power of $(\mu_n)_S$ over the constant sheaf of rings $\ZZ_S$, with $(\mu_n^{\otimes0})_S=(\ZZ/n)_S$. If $n$ is invertible on $S$, this is locally constant cyclic of order $n$. The usual tensor identities hold, and the construction is compatible with pullback.

For a codimension-$1$ regular pair $(Y,X)$, suppose locally that $Y$ is defined by $x=0$. There is an exact sequence
\[
0\to(\GG_m)_X\to j_*(\GG_m)_U\to \ZZ_Y\to0.
\]
Applying Lemma 1.3 gives canonical isomorphisms
\begin{equation}
   (R^1j_*)\mu_n\simeq
   j_*(\GG_m)_U/(j_*(\GG_m)_U)^n\simeq(\ZZ/n)_Y, \tag{3.3.1}
\end{equation}
independent of the chosen equation $x$. Consequently there is a canonical isomorphism
\begin{equation}
   \varphi:\ZZ/n\xrightarrow{\sim}(R^1i^!)\mu_n
   \simeq i^*(R^1j_*)\mu_n. \tag{3.3.2}
\end{equation}

\begin{statement}{Theorem 3.4}
Assuming resolution of singularities, let $C$ be the family of regular pairs of excellent schemes of equal characteristic. There exists a unique rule assigning to every regular pair $(Y,X)$ of codimension $c$ a canonical isomorphism of sheaves
\[
   \varphi(Y,X):\ZZ/n\xrightarrow{\sim}
   \mathcal H_Y^{2c}(X,\mu_n^{\otimes c})=(R^{2c}i^!)\mu_n^{\otimes c}
\]
satisfying the following conditions.
\begin{enumerate}[label=(\roman*)]
\item For $c=1$, if $Y$ is defined in $X$ by $x=0$, the isomorphism is (3.3.2).
\item It is compatible with morphisms of regular pairs: if $(Y',X')\to(Y,X)$ is such a morphism and $g:Y'\to Y$ is induced, the pullback of $\varphi(Y,X)$ agrees with $\varphi(Y',X')$.
\item It is transitive. If $(Z,Y,X)$ is a regular triple, with codimensions $a=\codim(Y,X)$, $b=\codim(Z,Y)$, and $c=a+b=\codim(Z,X)$, then the isomorphism for $(Z,X)$ is obtained from the isomorphism for $(Z,Y)$ followed by the isomorphism for $(Y,X)$ through the spectral sequence
\[
   (R^pv^!)(R^qu^!)F\Longrightarrow(R^{p+q}w^!)F.
\]
In other words the standard square comparing
\[
   \varphi(Z,X)
   \quad\text{and}\quad
   \xi\,(R^{2b}v^!)(\varphi(Y,X)\otimes \mu_n^{\otimes b})\,\varphi(Z,Y)
\]
is commutative.
\end{enumerate}
\end{statement}

\begin{prooftext}
Existence and uniqueness are proved by induction on $c$. For $c=1$, define $\varphi$ locally by (3.3.1). Since the construction is independent of the local parameter, it glues. Compatibility with morphisms follows directly from the definition and Lemma 1.3; transitivity is void in codimension $1$.

Assume that the result is known in codimension $<c$, with $c>1$. In a regular triple $(Z,Y,X)$, all the already defined arrows appearing in the transitivity diagram are compatible with morphisms of regular triples. Put
\[
   \Psi(Z,Y,X)=\xi\,\eta\,\varphi(Z,Y),
\]
where $\eta$ is the map obtained by applying $R^{2b}v^!$ to $\varphi(Y,X)$ after the Tate twist, and $\xi$ is the edge isomorphism supplied by the spectral sequence. This $\Psi$ is compatible with morphisms of triples.

Given a regular pair $(Z,X)$ of codimension $c>1$, one can locally choose a regular triple $(Z,Y,X)$ with $0<\codim(Y,X)<c$. Everything is therefore reduced to showing that $\Psi(Z,Y,X)$ is independent of the intermediate $Y$. A further transitivity chase reduces the question to the case where $(Y,X)$ has codimension $1$.

Let $(Z,Y_0,X)$ and $(Z,Y_1,X)$ be two such triples. Since the assertion may be checked fibrewise, assume $X$ strictly local and let $Y_i$ be defined by $f_i=0$. Put $\bar X=\Spec \Oo_X[t]$, let $\bar Z=Z\times_X\bar X$, and define $\bar Y\subset\bar X$ by
\[
   (f_1-f_0)t+f_0=0.
\]
This interpolates between $Y_0$ and $Y_1$ along the sections $t=0$ and $t=1$. Removing from $\bar X$ the closed set where $\bar Y$ is not regular, we obtain a regular triple $(\bar Z,\bar Y,\bar X)$ and two morphisms of triples
\[
   (Z,Y_i,X)\longrightarrow(\bar Z,\bar Y,\bar X),\qquad i=0,1.
\]
Let
\[
   F=(R^{2c}w^!)\mu_n^{\otimes c}
\]
on $Z$. By Theorem 3.2 it is locally constant and isomorphic to $\ZZ/n$. The analogous sheaf $\bar F$ on $\bar Z$ is the pullback of $F$. The fibres of $\bar Z\to Z$ are non-empty and connected, and the morphism is smooth and open; therefore
\[
   H^0(Z,F)\longrightarrow H^0(\bar Z,\bar F)
\]
is bijective. Pullback by the two sections $t=0$ and $t=1$ gives the same inverse. The global sections corresponding to the two isomorphisms $\Psi(Z,Y_i,X)$ are thus equal, because both are the pullbacks of the section corresponding to $\Psi(\bar Z,\bar Y,\bar X)$. Hence $\Psi(Z,Y_0,X)=\Psi(Z,Y_1,X)$.

This proves independence of the intermediate $Y$, hence uniqueness in codimension $c$, and existence follows by gluing. \qedtext
\end{prooftext}



\section*{4. Local acyclicity of a regular morphism}

\begin{statement}{Theorem 4.1 (local acyclicity)}
Assuming resolution of singularities, let
\[
   g:X'\longrightarrow X
\]
be a regular morphism of excellent schemes of equal characteristic. Then $g$ is universally locally acyclic for torsion coefficients of order prime to the characteristic.
\end{statement}

\begin{prooftext}
The hypothesis is stable under finite type base change. By the usual criterion for universal local acyclicity, it is enough to prove local acyclicity for $n$. Applying the criterion of Expos\'e XV, 1.17, take a geometric point $\bar x'$ of $X'$ and its image $\bar x$ in $X$. Strict localization preserves regular morphisms, hence we reduce to the case
\[
   X=\Spec A,
   \qquad X'=\Spec A',
\]
with $A$ and $A'$ strictly local, and with $g$ local. By the criterion just cited, it remains to show that the geometric fibres of $g$ are acyclic for $n$.

After replacing $X$ by an integral closed subscheme and $X'$ by the corresponding closed subscheme, it is enough to consider the generic geometric fibre $X'_{\bar x}$, where $\bar x=\Spec\bar K$ and $\bar K$ is the separable closure of the rational function field of $X$. Write
\[
   \bar x=\varprojlim_\alpha U_\alpha
\]
with $U_\alpha$ ranging over a filtered category of connected etale $X$-schemes, which may be taken regular. Let $B_\alpha$ be the normalization of $A$ in the rational function field of $U_\alpha$, put
\[
   U'_
\alpha=X'\times_XU_\alpha,
   \qquad B'_
\alpha=A'\otimes_AB_\alpha.
\]
Then $U_\alpha$ is a regular open subscheme of $\Spec B_\alpha$, and $B_\alpha\to B'_
\alpha$ is regular. Lemma 2.4 (or the comparison established there) gives
\[
   H^q(U_\alpha,\mu_n)\xrightarrow{\sim}H^q(U'_
\alpha,\mu_n)
\]
for every $q$. Since $X'_{\bar x}$ is the projective limit of the $U'_
\alpha$,
\begin{align*}
H^q(X'_{\bar x},\mu_n)
&=\varinjlim_
\alpha H^q(U'_
\alpha,\mu_n)\\
&=\varinjlim_
\alpha H^q(U_\alpha,\mu_n)
 =H^q(\bar x,\mu_n)=0\qquad(q>0).
\end{align*}
This proves local acyclicity, hence universal local acyclicity. \qedtext
\end{prooftext}

\begin{statement}{Corollary 4.2 (base change by a regular morphism)}
Let
\[
\xymatrix@C=2.0cm@R=1.0cm{
X\ar[d]^f&X'\ar[l]_{g'}\ar[d]^{f'}\\
Y&Y'\ar[l]_g
}
\]
be Cartesian, with $g:Y'\to Y$ a regular morphism of excellent preschemes of equal characteristic, and with $f$ quasi-compact and quasi-separated. For every abelian torsion sheaf $F$ on $X$ of order prime to the characteristic, the base-change morphisms
\[
   g^*R^qf_*F\longrightarrow R^qf'_*{g'}^*F
\]
are isomorphisms for all $q\ge0$.
\end{statement}

This follows from Theorem 4.1 and the local acyclicity theorem of Expos\'e XVI, 1.1.

\begin{statement}{Corollary 4.3}
Let $A$ be an excellent strictly local ring of equal characteristic, let $f:X\to\Spec A$ be a morphism, and let $A\to A'$ be a regular morphism with $A'$ also excellent, strictly local, and of equal characteristic. Denote by a prime the base change along $A\to A'$. For every torsion sheaf $F$ on $X$ of order prime to the characteristic of $A$,
\[
   H^q(X,F)\xrightarrow{\sim}H^q(X',F')
\]
for all $q\ge0$.
\end{statement}

In particular this applies when $A'$ is the completion $\widehat A$ of $A$. In that case the hypothesis that $A$ be strictly local may be weakened to the hypothesis that $A$ be local henselian.

\section*{5. Finiteness theorem}

\begin{statement}{Theorem 5.1 (finiteness)}
Assuming resolution of singularities, let $f:X\to S$ be a morphism of finite type between excellent schemes of equal characteristic, and let $F$ be a constructible abelian torsion sheaf on $X$ of order prime to the characteristic. Then all the sheaves $R^qf_*F$ are constructible.
\end{statement}

This generalizes Expos\'e XVI, 5.1. The proof is essentially the same: one uses the purity theorem 3.2 in place of Expos\'e XVI, 3.7.

\begin{statement}{Corollary 5.2}
Let $A$ be an excellent strictly local ring of equal characteristic, let $f:X\to\Spec A$ be of finite type, and let $F$ be a constructible abelian torsion sheaf on $X$ of order prime to the characteristic of $A$. Then all groups $H^q(X,F)$ are finite.
\end{statement}

\section*{6. Cohomological dimension of affine morphisms}

Let $f:X\to Y$ be a morphism of finite type between excellent preschemes, and let $F$ be a sheaf on $X$. Recall that, when $Y$ is local, Expos\'e XIV defined an integer $\delta(F)$. For arbitrary $Y$ we view it as a function by setting
\begin{equation}
   \delta(F,f,y)=\delta(F'), \tag{6.0}
\end{equation}
where $Y'=\Spec\Oo_{Y,y}$, $X'=X\times_YY'$, and $F'$ is the inverse image of $F$ on $X'$.

\begin{statement}{Theorem 6.1}
Assuming resolution of singularities, let $f:X\to Y$ be an affine morphism of finite type, with $Y$ an excellent scheme of equal characteristic, and let $F$ be an abelian torsion sheaf on $X$. Then $R^qf_*F$ is zero at every point $y\in Y$ for which $\delta(F,f,y)<q$.
\end{statement}

Equivalently, in local form:

\begin{statement}{Theorem 6.1 bis}
Let $A$ be an excellent strictly local ring of equal characteristic, and let
\[
   f:X\longrightarrow Y=\Spec A
\]
be affine of finite type. If $F$ is a torsion sheaf on $X$, then
\[
   H^q(X,F)=0\qquad(q>\delta(F)).
\]
\end{statement}

For torsion of characteristic $p$ when $p$ is the characteristic, such statements have no content in this form; hence one may suppose the torsion prime to the characteristic.

\begin{statement}{Corollary 6.2}
Let $A$ be a henselian excellent ring of equal characteristic with residue field $k$, and let $U\subset\Spec A$ be affine. For a prime $\ell$,
\[
   \operatorname{cd}_\ell U\leq \operatorname{cd}_\ell k+\dim A.
\]
\end{statement}

Passing to the strict localization $\widetilde A$ reduces the assertion by the Hochschild-Serre spectral sequence to the case where $A$ is strictly local, where it is a special case of Theorem 6.1 bis.

\begin{statement}{Corollary 6.3}
Let $A$ be henselian, excellent, and of equal characteristic, with residue field $k$, and let $R(A)$ be the field of rational functions on $\Spec A$. For a prime $\ell$,
\[
   \operatorname{cd}_\ell R(A)\leq \operatorname{cd}_\ell k+\dim A.
\]
If $\ell$ is invertible in $k$, and either $\ell\ne2$ or $k$ is not orderable, equality holds.
\end{statement}

The equality follows from the corresponding lower-bound theorem in Expos\'e X.

\begin{statement}{Corollary 6.4}
Let $X$ be a noetherian excellent scheme of equal characteristic, let $R(X)$ be its ring of rational functions, and let $\ell$ be invertible on $X$. Suppose moreover that $\ell\ne2$, or that no residue field of $X$ is orderable. Then
\[
   \operatorname{cd}_\ell(X)\leq \operatorname{cd}_\ell R(X)+\dim X
      \leq 2\operatorname{cd}_\ell R(X).
\]
\end{statement}

\paragraph{Proof of Theorem 6.1.}
The proof begins with the reduction of Expos\'e XIV, Section 3. Induct on $\delta$. The assertion for $\delta\le d$ is equivalent to the local assertion 6.1 bis for $\delta\le d$. The case $\delta=0$ is trivial. When $\delta\le1$, one may suppose $F$ constructible, since a torsion sheaf is the filtered union of its constructible subsheaves. The support of $F$ is then contained in a closed subset $Z$ whose points have $\delta\le1$. Each reduced irreducible component of $Z$ is either affine of finite type of dimension $\le1$ over the closed point of $Y$, or quasi-finite over a one-dimensional closed subscheme of $Y$. In every case its cohomological dimension is at most $1$. A standard exact sequence decomposing $F$ over the irreducible components of $Z$ gives the result.

Assume now that the theorem is known for values of $\delta<d$, and prove it for $\delta(F)\le d$, with $d\ge2$.

\begin{statement}{Lemma 6.5}
It is enough to treat the case where $Y$ is strictly local and
\[
   X=\mathbb A^1_Y=\Spec\Oo_Y[t].
\]
\end{statement}

The proof is the same as Expos\'e XIV, 4.2. One first reduces to $Y$ strictly local. Embedding $X$ in affine space over $Y$, induction on the affine dimension reduces the problem from $\mathbb A^N_Y$ to $\mathbb A^1_{\mathbb A^{N-1}_Y}$. The Leray spectral sequence for the projection $g:\mathbb A^N_Y\to\mathbb A^{N-1}_Y$ shows that it is enough to establish
\[
   \delta(R^qg_*F)\le d-q.
\]
This follows from the inequality
\[
   \delta(F,g,z)+\delta(z)\le \delta(F),
\]
whose verification is formal from the definition of $\delta$.

\begin{statement}{Lemma 6.6}
To prove Theorem 6.1 for $\delta\le d$, it is enough to prove the following assertion. Let $B$ be a complete normal local ring of equal characteristic and dimension $d$. Let $U\subset\Spec B$ be an affine open, and suppose that some $x\in\rad B$ is invertible on $U$. Then
\[
   H^q(U,\ZZ/\ell)=0\qquad(q>d).
\]
\end{statement}

This lemma is itself a special case of Corollary 6.2, but it is used here as the key reduction. Starting from Lemma 6.5, compactify $\mathbb A^1_Y$ into $\mathbb P^1_Y$ and use the spectral sequence for the open immersion $i:\mathbb A^1_Y\hookrightarrow\mathbb P^1_Y$. The higher $R^qi_*F$ are supported on the section at infinity, and proper cohomological dimension for $\mathbb P^1_Y$ leaves only the stalks at infinity to be controlled. These stalks are cohomology groups of strict localizations of closed subschemes of $\mathbb P^1_Y$ with one element inverted. Normalizing, and then completing, reduces by Corollary 4.3 to the situation of Lemma 6.6. The induction hypothesis handles the cases of dimension $<d$.

We next prove Lemma 6.6 in characteristic zero, and then finish the positive-characteristic case in Section 7.

\paragraph{6.7. Formal separability.}
Let $A$ be an equidimensional complete local $k$-algebra of equal characteristic, with $A/\rad A$ finite over $k$. We say that $A$ is \emph{formally separable} over $k$ if it contains a subalgebra $k[[x_1,\ldots,x_n]]$ such that
\[
   \Spec A\longrightarrow \Spec k[[x_1,\ldots,x_n]]
\]
is finite and generically etale.

\begin{statement}{Lemma 6.7}
Let $k\subset B$ be a field over which $B/\rad B$ is finite. Under the induction hypothesis on $d$, the conclusion of Lemma 6.6 holds if the hypotheses of 6.6 hold and, in addition, one of the following conditions is satisfied:
\begin{enumerate}[label=(\roman*)]
\item $k$ has characteristic zero;
\item there exists an element $x\in\rad B$, invertible on $U$, such that $B/(x)$ is a formally separable $k$-algebra.
\end{enumerate}
\end{statement}

\begin{prooftext}
Choose parameters $x_1=x,x_2,\ldots,x_d$ in $\rad B$ generating an ideal primary for $\rad B$. Then $B$ is finite over
\[
   A=k[[x_1,\ldots,x_d]].
\]
In case (ii) one can choose $x_2,\ldots,x_d$ so that $B/x_1B$ is generically etale over $k[[x_2,\ldots,x_d]]$. Hence $B$ is generically etale over $A$; this is automatic in characteristic zero.

Choose a primitive element $b$ for the extension of fraction fields and let $\varepsilon\in k[[x_1,
\ldots,x_d]]$ be the discriminant of its monic irreducible equation. In case (ii) $b$ may be chosen so that $x_1$ does not divide $\varepsilon$. Write
\[
   \Spec B-U=V(x_1)\cup Y,
\]
where $Y$ contains no irreducible component of $V(x_1)$. After a suitable change of the remaining parameters, one arranges that $V(x_d)$ is not contained in the image of $Y$, and in case (ii) is not contained in $V(\varepsilon)$.

\begin{statement}{Lemma 6.8}
Let
\[
   A^0=k[[x_1,\ldots,x_{d-1}]]\{x_d\}
\]
be the strict localization of $k[[x_1,\ldots,x_{d-1}]][x_d]$ at the origin. Under the preceding choices there exist an $A^0$-algebra $B^0$, an affine open $U^0\subset\Spec B^0$, and an isomorphism
\[
   A\otimes_{A^0}B^0\simeq B
\]
such that the open of $\Spec B$ induced from $U^0$ is $U$.
\end{statement}

The proof is parallel to Lemma 1.5. One algebraizes the finite algebras defining $B$ and the closed complement $Y$ from $A$ to $A^0$. In characteristic zero this follows from Artin approximation; in case (ii) it follows from the same Weierstrass argument as before, because of the generically etale condition. Then set
\[
   U^0=\Spec B^0-(\Spec C^0\cup V(x_1)).
\]
The base change to $A$ is exactly $U$.

Applying Corollary 4.3 once more gives
\[
   H^q(U,\ZZ/\ell)\simeq H^q(U^0,\ZZ/\ell).
\]
Now $\Spec A^0$ is a limit of affine schemes etale over $k[[x_1,\ldots,x_{d-1}]][x_d]$. Thus $U^0$ is a projective limit of affine schemes $U^0_\alpha$ of finite type over
\[
   V^0=\Spec k[[x_1,\ldots,x_{d-1}]][1/x_1],
\]
whose fibres over $V^0$ have dimension at most $1$. For the structural morphism $f^0_\alpha:U^0_\alpha\to V^0$ one has, at $v\in V^0$,
\[
   \delta(\ZZ/\ell,f^0_\alpha,v)\le \dim\Oo_{V^0,v}+1\le d-1.
\]
The induction hypothesis applied to $f^0_\alpha$, and then to the open immersion $V^0\hookrightarrow\Spec k[[x_1,\ldots,x_{d-1}]]$, gives
\[
   H^p(V^0,R^qf^0_{\alpha*}\ZZ/\ell)=0\qquad(p+q>d).
\]
The Leray spectral sequence therefore gives $H^n(U^0_\alpha,\ZZ/\ell)=0$ for $n>d$. Passing to the limit gives $H^n(U^0,\ZZ/\ell)=0$ for $n>d$. This proves Lemma 6.7, hence Theorem 6.1 in characteristic zero. \qedtext
\end{prooftext}

\section*{7. Affine morphisms -- end of the proof}

It remains, in characteristic $p>0$ and still under the resolution hypothesis, to finish Lemma 6.6. We have reduced to the following statement for $d\ge2$.

\begin{statement}{Lemma 7.1}
Let $B$ be a complete normal local ring of equal characteristic, with separably closed residue field, and of dimension $d$. Let $U\subset\Spec B$ be an affine open such that some $x\in\rad B$ is invertible on $U$. Then
\[
   H^q(U,\ZZ/\ell)=0\qquad(q>d).
\]
Moreover, by induction, Theorem 6.1 may be assumed known for $\delta<d$, and Lemma 7.1 may be assumed known in the special case 6.7(ii).
\end{statement}

Normality is inessential. If $\bar B$ is the normalization and $\bar U$ is the inverse image of $U$, then
\[
0\to(\ZZ/\ell)_U\to \pi_*(\ZZ/\ell)_{\bar U}\to C\to0
\]
has $\dim\operatorname{Supp}C<d$. By induction $H^q(U,C)=0$ for $q>d-1$, and $H^q(U,\pi_*\cdot)=H^q(\bar U,\cdot)$; hence
\begin{equation}
H^q(U,\ZZ/\ell)=0\ (q>d)
\quad\Longleftrightarrow\quad
H^q(\bar U,\ZZ/\ell)=0\ (q>d). \tag{7.2}
\end{equation}

\paragraph{7.3. Formal separability in characteristic \texorpdfstring{$p$}{p}.}
We recall without proof the facts needed below. Let $A$ be a complete noetherian local $k$-algebra with residue field finite over $k$, and let $k'=k^{1/p}$. If $A'$ is the completion of the local ring of $A\otimes_kk'$, then $A\to A'$ is radicial, possibly infinite. This replacement does not change etale cohomology. The condition that an equidimensional algebra $A$ be formally separable over $k$ is expressed through the completed module of differentials $\widehat\Omega^1_{A/k}$, or through a Jacobian criterion. If $A$ is not analytically separable, then the corresponding $k'$-algebra $A'$ is not reduced.

\begin{statement}{Lemma 7.3}
Let $f(x_1,\ldots,x_n)\in k[[x_1,\ldots,x_n]]$. Suppose that for each $i$ there exists a unit $u_i$ such that $u_if$ is a series in $x_i^p$ and in the other variables. Then there exists a unit $v$ such that
\[
   vf\in k[[x_1^p,\ldots,x_n^p]].
\]
\end{statement}

Iterating the operation $(A,k)\mapsto(\bar A,k^{1/p})$, where $\bar A$ is the normalization after base change, one reduces to a formally separable situation. The same procedure may be applied to $(A/(x))_{\red}$ for a fixed $x\in\rad A$.

\begin{statement}{Corollary 7.4}
It is enough to prove Lemma 7.1 under the additional hypothesis that $B$ contains a field $k$ over which $B/\rad B$ is finite, and such that $(B/(x))_{\red}$ is a formally separable $k$-algebra.
\end{statement}

Now assume these conditions. Write
\[
   D(x)=\sum_i r_iX_i
\]
for the divisor of $x$, with $X_i$ irreducible closed subschemes of codimension $1$. Choose $y\in\rad B$ such that
\[
   D(y)=\sum_i(r_i-1)X_i+\Gamma,
\]
where $\Gamma$ contains none of the $X_i$. Blow up the ideal $(x,y)$:
\[
   f:Z\longrightarrow\Spec B.
\]
The two affine opens $\Spec B[y/x]$ and $\Spec B[x/y]$ cover $Z$. Let $B_1$ be the strict localization of $B[x/y]$ at the maximal ideal generated by $(\rad B,x/y)$. For each closed point $Q$ of the closed fibre of $\Spec B[y/x]\to\Spec B$, let $A_Q$ be the strict localization of $B[y/x]$ at $Q$, and let $U_Q$ be the inverse image of $U$ in $\Spec A_Q$.

\begin{statement}{Lemma 7.5}
To prove Lemma 7.1 for $B$, it is enough to prove, for $q>d$,
\begin{equation}
   H^q(\Spec B_1[y/x],\ZZ/\ell)=0, \tag{i}
\end{equation}
and, for every such closed point $Q$,
\begin{equation}
   H^q(U_Q,\ZZ/\ell)=0. \tag{ii}
\end{equation}
\end{statement}

\begin{prooftext}
There are open immersions
\[
   U\hookrightarrow\Spec B[1/x]\hookrightarrow\Spec B[y/x].
\]
Use the Leray spectral sequence for the resulting immersion $i:U\to\Spec B[y/x]$. By (ii), the stalks $(R^qi_*\ZZ/\ell)_Q$ vanish for $q>d$ at the closed points $Q$ of the closed fibre. At all other points $P$ one has $\delta(\ZZ/\ell,i,P)<d$; by the induction hypothesis, $R^qi_*\ZZ/\ell$ has support of $\delta$ at most $d-q$. Applying the induction hypothesis in the form 6.1 bis to the morphism $\Spec B[y/x]\to\Spec B$, one obtains
\[
   H^p(\Spec B[y/x],R^qi_*\ZZ/\ell)=0
\]
for $q>0$ and $p>d-q$. Thus the proof of Lemma 7.1 is reduced to
\[
   H^p(\Spec B[y/x],i_*\ZZ/\ell)=0\qquad(p>d).
\]
The exact sequence
\begin{equation}
   0\to \ZZ/\ell\to i_*(\ZZ/\ell)_U\to C\to0 \tag{7.6}
\end{equation}
with $\delta(C)<d$, together with induction, reduces further to
\begin{equation}
   H^p(\Spec B[y/x],\ZZ/\ell)=0\qquad(p>d). \tag{7.7}
\end{equation}

Let $j:\Spec B[y/x]\hookrightarrow Z$ be the open immersion; the complement is the divisor at infinity, which in the other chart $\Spec B[x/y]$ is $V(x/y)$. The sheaves $R^qj_*\ZZ/\ell$ for $q>0$ are concentrated on this closed set, which maps by a closed immersion to $\Spec B$. Hence their higher cohomology on $Z$ vanishes. Since $Z\to\Spec B$ is proper with one-dimensional closed fibre, proper cohomological dimension gives $H^p(Z,j_*\ZZ/\ell)=0$ for $p>2$. The Leray spectral sequence for $j$ identifies, for $q>2$, the groups $H^q(\Spec B[y/x],\ZZ/\ell)$ with the stalks at the point at infinity. These stalks are cohomology groups of the form
\[
   H^q(\Spec B_1[y/x],\ZZ/\ell).
\]
Thus (i), plus the induction already used, implies (7.7), and the lemma follows. \qedtext
\end{prooftext}

\begin{prooftext}[Proof of Lemma 7.5(i)]
Let $C=V(x/y)\subset\Spec B[x/y]$ with its reduced structure. The morphism $Z\to\Spec B$ identifies $C$ with the closed subscheme $X=\bigcup_iX_i$. Since
\[
   D(x)=\sum r_iX_i,
   \qquad D(y)=\sum(r_i-1)X_i+\Gamma,
\]
the element $x/y$ generates the ideal of $X_i$ at the generic point of $X_i$; there the local ring is a discrete valuation ring because $B$ is normal. Hence $\Spec B[x/y]\to\Spec B$ is an isomorphism near these generic points, and $x/y$ vanishes to order $1$ along each irreducible component of $C$.

The same remains true after passing to the strict localization $B_1$. Thus $B_1/(x/y)$ is a complete $k$-algebra which is formally separable. Corollary 4.3 allows one to replace $B_1$ by its completion; then (7.2) allows one to replace it by its normalization. Since the normalization is still an isomorphism at the generic points of the components of $C$, the quotient by $x/y$ remains formally separable. This is precisely the special case 6.7(ii), so (i) follows.
\end{prooftext}

\begin{prooftext}[Proof of Lemma 7.5(ii)]
Fix a closed point $Q$ of the closed fibre of $\Spec B[y/x]\to\Spec B$. The strict localization $A_Q$ is a filtered limit of etale finite-type $B[y/x]$-algebras $A$. Let $U_A$ be the inverse image of $U$ in $\Spec A$. Since $x$ is invertible on $U_A$, the morphism $U_A\to\Spec B$ is etale, and $U_A$ is normal. Let $C$ be the normalization of $B$ in the rational function field of $A$. Then $U_A$ is an open subscheme of $\Spec C$.

It is enough to show that, for each such $A$, the image of $H^q(U_A,\ZZ/\ell)$ in $H^q(U_Q,\ZZ/\ell)$ is zero for $q>d$. Since $x$ is invertible on $U_Q$, the morphism $U_Q\to\Spec A$ factors through $\Spec A[1/f]$ for any element $f\in B$ divisible by $x$ in $B[y/x]$ and such that $f/x$ is nonzero at $Q$. Thus it suffices to find a suitable such $f$ for which
\[
   H^q(U_{A,f},\ZZ/\ell)=0\qquad(q>d),
\]
where $U_{A,f}=U_A-V(f)$.

Choose $c\in k$ such that $y/x+c$ does not vanish at $Q$, and let $J$ be the ideal of the closed set $\Gamma\cap X$ in $B$. If
\[
   f\equiv y+cx\pmod{J^N},\qquad N>0, \tag{7.8}
\]
then $f$ is divisible by $x$ in $B[y/x]$ and $f/x\equiv y/x+c\pmod J$, hence $f/x$ is nonzero at $Q$. Since $\Gamma\cap X$ has codimension $2$, one can choose such an element, call it $z$, so that $V(z)\cap V(y)$ has codimension $2$. Choose $x_3,\ldots,x_d$ so that
\[
   y,z,x_3,\ldots,x_d
\]
is a system of parameters of $B$. Then $B$, and also $C$, is finite over $k[[y,z,x_3,\ldots,x_d]]$.

Let $C^0$ be the normalization of this power-series ring in the separable closure of $k((y,z,x_3,\ldots,x_d))$ inside the rational function field of $C$. Then $C$ is radicial over $C^0$, and $C^0$ is generically etale over the power-series ring. Choose a linear combination
\[
   f=az+by\qquad(a,b\in k)
\]
such that $\Spec C^0\to\Spec k[[z,y,x_3,
\ldots,x_d]]$ is etale above the generic point of $V(f)$, and such that $f/x$ is nonzero at $Q$ in $\Spec B[y/x]$. By radicial invariance of etale cohomology, it is enough to prove
\[
   H^q(U^0_{A,f},\ZZ/\ell)=0\qquad(q>d),
\]
where $U^0_{A,f}$ is the corresponding affine open of $\Spec C^0$. But $f$ is invertible on this open and $C^0/(f)$ is formally separable over $k$. This is again the special case 6.7(ii). This proves (ii), hence Lemma 7.5, Lemma 7.1, and Theorem 6.1. \qedtext
\end{prooftext}

\subsection*{Bibliography of Expos\'e XIX}

\begin{itemize}
\item S. Abhyankar, \emph{Local uniformization on algebraic surfaces over ground fields of characteristic $p\ne0$}, Annals of Mathematics 63 (1963), and \emph{Resolution of singularities of embedded algebraic surfaces}, Academic Press, 1966.
\item S. Abhyankar, \emph{Resolution of singularities of arithmetical surfaces}, in \emph{Arithmetical Algebraic Geometry}, New York, 1966.
\item M. Artin, \emph{Etale coverings of schemes over hensel rings}, American Journal of Mathematics 88, no. 4, 915--934 (1966).
\item H. Hironaka, \emph{Resolution of singularities of an algebraic variety over a field of characteristic zero I, II}, Annals of Mathematics 79 (1964).
\item K. Weierstrass, \emph{Einige auf die analytischen Functionen mehrerer Ver\"anderlichen sich beziehende S\"atze}, Mathematische Werke von K. Weierstrass, vol. 2, Berlin, 1895.
\item N. Bourbaki, \emph{Commutative Algebra}, Chapter 3.
\end{itemize}

\bigskip

\end{document}
